\begin{figure*}[t]
\centering
\scriptsize
\setlength{\fboxsep}{1em}
\setlength{\fboxrule}{1pt}
\fcolorbox{black}{white}{
\parbox{0.955\textwidth}{
\begin{alltt}


You will be provided instructions for completing a task followed by a task to complete.

Instructions:

1. Read the question carefully. The dataset contains multiple-choice questions related to college-level chemistry. 

2. Understand the context of the question. The topics can range from atomic structure, chemical reactions, stoichiometry, thermodynamics, chemical bonding, to organic chemistry. 

3. Evaluate each choice. The dataset provides a list of choices for each question. 

4. Identify the correct answer. The correct answer is indicated by a class label in the dataset. 

5. Avoid the negative answer. The dataset also provides a negative answer, which is an incorrect option. 

6. Use your knowledge of college-level chemistry to answer the question. You might need to recall facts, apply formulas, or use logical reasoning based on the principles of chemistry. 

7. If the question is complex, break it down into smaller parts and address each part separately. 

8. Always provide an explanation for your answer. This helps to demonstrate your understanding of the topic. 

9. Review your answer before finalizing it. Ensure that it is accurate and complete. 

10. Keep learning and updating your knowledge about college-level chemistry. The field is vast and the questions can cover any topic within it.

Question: Which of the following is true for Br2 at standard temperature and pressure?

A. It is a colorless gas.

B. It is a red-brown volatile liquid.

C. It is a colorless volatile liquid.

D. It is a yellow metallic solid.

Follow the instructions to generate an explanation that reasons towards the correct answer to the task above. End the explanation with the correct answer. [/INST] Explanation: 

The question asks about the state of Br2 at standard temperature and pressure. The options are:

A. Colorless gas

B. Red-brown volatile liquid

C. Colorless volatile liquid

D. Yellow metallic solid

To answer this question, we need to recall the properties of bromine (Br2). Bromine is a reddish-brown liquid at room temperature and pressure. It has a density of 2.96 g/mL and a boiling point of 58.8°C.

Therefore, the correct answer is B. Red-brown volatile liquid.

Explanation:

The statement A is incorrect because Br2 is not a gas at standard temperature and pressure. It is a liquid.

The statement C is incorrect because Br2 is not a colorless liquid. It has a reddish-brown color.

The statement D is incorrect because Br2 is not a solid, and it does not have a yellow color.

Therefore, the correct answer is B. Red-brown volatile liquid.

Therefore, the correct multiple choice label (just the letter) to the task is below.

Answer: B
\end{alltt}
}}
\caption{A complete running example of MMLU (College Chemistry).}
\label{fig:mmlu_subject=college_chemistry_example}
\end{figure*}
